\documentclass[../DoAn.tex]{subfiles}
\begin{document}

\section{Thiết kế kiến trúc}
\label{section:4.1}

\subsection{Lựa chọn kiến trúc phần mềm}
\label{subsection:4.1.1}

Hệ thống Aikari được xây dựng dựa trên kiến trúc ba lớp (Three-Tier Architecture) kết hợp với mô hình Client-Server, phù hợp với yêu cầu của một ứng dụng web hiện đại cần xử lý nhiều người dùng đồng thời, tích hợp AI, và đảm bảo tính mở rộng trong tương lai.

Kiến trúc ba lớp chia ứng dụng thành ba tầng độc lập: tầng trình bày (Presentation Layer), tầng nghiệp vụ (Business Logic Layer), và tầng dữ liệu (Data Access Layer). Kiến trúc này có ưu điểm về tách biệt trách nhiệm, khả năng tái sử dụng, mở rộng và bảo mật.

Hệ thống Aikari áp dụng kiến trúc ba lớp với các thành phần cụ thể như sau:

\textbf{Tầng trình bày} được triển khai bằng React Router v7 và TypeScript, bao gồm các component React được tổ chức theo cấu trúc components (ui, shared, layout, pages, admin, dashboard), routes, và redux. State management sử dụng Redux Toolkit, routing quản lý bởi React Router v7, và styling sử dụng Tailwind CSS v4 với Radix UI components. Tầng trình bày giao tiếp với tầng nghiệp vụ thông qua REST API sử dụng HTTP/HTTPS.

\textbf{Tầng nghiệp vụ} được triển khai bằng FastAPI và Python, được tổ chức theo mô hình module hóa với các API routes (login, users, studysets, learning, tests, chatbot, classes, admin, session, categories, ai), services (learning\_service, chatbot\_service, generation\_service, dify\_service, analytics\_service, session\_service, admin\_service, email\_service), auth utilities, và schemas cho validation. Tầng này xử lý các logic phức tạp như tính toán thuật toán SM-2, quản lý conversation với AI, tích hợp Dify workflows, phân quyền, và tính toán thống kê.

\textbf{Tầng dữ liệu} được triển khai bằng SQLModel (ORM) và PostgreSQL, bao gồm các models (User, StudySet, Term, StudyActivity, ProgressSummary, Test, TestQuestion, TestAttempt, TestAnswer, Class, ClassMember, ClassStudySet, ChatConversation, ChatMessage, Category, AIGeneratedContents, LearningSession, SessionReview, RefreshToken, ReattemptRequest), database connection management, và Alembic migrations. SQLModel kết hợp SQLAlchemy và Pydantic, cho phép định nghĩa models một lần và sử dụng cho cả database schema và API validation.

Hệ thống có một số cải tiến so với kiến trúc ba lớp truyền thống: Service Layer Pattern để tách biệt API routes và services, Dependency Injection mạnh mẽ của FastAPI, hỗ trợ Async/Await, Type Safety với TypeScript và Pydantic, và tích hợp External Service (Dify) thông qua DifyService.

\subsection{Thiết kế tổng quan}
\label{subsection:4.1.2}

Hệ thống Aikari được tổ chức thành các gói (package) chính theo kiến trúc ba lớp, đảm bảo sự phụ thuộc rõ ràng và tuân thủ các nguyên tắc thiết kế. Biểu đồ gói UML mô tả cấu trúc tổng quan của hệ thống được trình bày trong Hình \ref{fig:package-diagram}.

\begin{figure}[H]
    \centering
    \includegraphics[width=0.9\textwidth]{Hinhve/package-diagram.png}
    \caption{Biểu đồ phụ thuộc gói của hệ thống Aikari}
    \label{fig:package-diagram}
\end{figure}

Các gói được tổ chức theo các tầng như sau:

\textbf{Tầng trình bày (Presentation Layer)}:
\begin{itemize}
    \item \texttt{frontend.app.components}: Chứa các React components (ui, shared, layout, pages, admin, dashboard)
    \item \texttt{frontend.app.redux}: Quản lý state toàn cục với Redux Toolkit
    \item \texttt{frontend.app.routes}: Định nghĩa routing và navigation
    \item \texttt{frontend.app.utils}: Các utilities và helper functions
\end{itemize}

\textbf{Tầng nghiệp vụ (Business Logic Layer)}:
\begin{itemize}
    \item \texttt{backend.app.api.routes}: Các API endpoints xử lý HTTP requests (login, users, studysets, learning, tests, chatbot, classes, admin, session, categories, ai)
    \item \texttt{backend.app.services}: Logic nghiệp vụ chính (learning, chatbot, generation, analytics, session, admin, email, dify)
    \item \texttt{backend.app.schemas}: Pydantic schemas cho validation
    \item \texttt{backend.app.api.deps}: Dependency injection
    \item \texttt{backend.app.core}: Cấu hình và utilities cốt lõi (config, security, db)
\end{itemize}

\textbf{Tầng dữ liệu (Data Access Layer)}:
\begin{itemize}
    \item \texttt{backend.app.models}: Định nghĩa các database models
    \item \texttt{backend.app.core.db}: Quản lý database connection và sessions
    \item \texttt{backend.app.alembic}: Database migrations
\end{itemize}

Các nguyên tắc thiết kế được tuân thủ:
\begin{itemize}
    \item Gói tầng dưới không phụ thuộc gói tầng trên
    \item Không phụ thuộc bỏ qua tầng
    \item Các gói trong cùng tầng không phụ thuộc lẫn nhau (trừ khi cần thiết)
    \item Tầng trình bày chỉ phụ thuộc tầng nghiệp vụ thông qua API
    \item Tầng nghiệp vụ chỉ phụ thuộc tầng dữ liệu thông qua models và database sessions
\end{itemize}

\subsection{Thiết kế chi tiết gói}
\label{subsection:4.1.3}

\subsubsection{Thiết kế gói Services}

Gói \texttt{backend.app.services} chứa logic nghiệp vụ chính của hệ thống. Biểu đồ thiết kế gói này được trình bày trong Hình \ref{fig:services-package}.

\begin{figure}[H]
    \centering
    \includegraphics[width=0.9\textwidth]{Hinhve/services-package.png}
    \caption{Thiết kế gói Services}
    \label{fig:services-package}
\end{figure}

Các lớp chính trong gói Services:
\begin{itemize}
    \item \texttt{LearningService}: Triển khai thuật toán SM-2, quản lý học tập và tiến độ
    \item \texttt{SessionService}: Quản lý learning sessions và session-based reviews
    \item \texttt{ChatbotService}: Xử lý conversation với AI, quản lý state và intent
    \item \texttt{GenerationService}: Sinh nội dung học tập (test, quiz, paragraph) từ AI
    \item \texttt{DifyService}: Tích hợp với Dify workflows để gọi AI
    \item \texttt{AnalyticsService}: Tính toán thống kê và phân tích dữ liệu học tập
    \item \texttt{AdminService}: Thống kê và quản trị hệ thống
    \item \texttt{EmailService}: Gửi email và quản lý email templates
\end{itemize}

Ngoài ra, hệ thống còn có các utility functions cho authentication (generate\_password\_reset\_token, verify\_password\_reset\_token) trong module \texttt{auth\_service}.

Các quan hệ giữa các lớp:
\begin{itemize}
    \item \texttt{LearningService} phụ thuộc vào \texttt{SpacedRepetitionSM2} (dependency) để tính toán thuật toán SM-2
    \item \texttt{ChatbotService} và \texttt{GenerationService} phụ thuộc vào \texttt{DifyService} (dependency) để gọi AI
    \item Tất cả services phụ thuộc vào các models (User, StudySet, Term, etc.) để truy xuất và cập nhật dữ liệu
\end{itemize}

\subsubsection{Thiết kế gói Models}

Gói \texttt{backend.app.models} định nghĩa các database models và quan hệ giữa chúng. Biểu đồ thiết kế gói này được trình bày trong Hình \ref{fig:models-package}.

\begin{figure}[H]
    \centering
    \includegraphics[width=0.9\textwidth]{Hinhve/models-package.png}
    \caption{Thiết kế gói Models}
    \label{fig:models-package}
\end{figure}

Các lớp chính trong gói Models:
\begin{itemize}
    \item \texttt{User}: Người dùng hệ thống
    \item \texttt{StudySet}: Bộ flashcard
    \item \texttt{Term}: Thẻ flashcard
    \item \texttt{StudyActivity}: Ghi nhận mỗi lần học
    \item \texttt{ProgressSummary}: Tổng hợp tiến độ học tập
    \item \texttt{LearningSession}: Quản lý learning sessions
    \item \texttt{SessionReview}: Ghi nhận review trong session
    \item \texttt{Test}: Bài kiểm tra
    \item \texttt{TestQuestion}: Câu hỏi trong bài kiểm tra
    \item \texttt{TestAttempt}: Lần làm bài kiểm tra
    \item \texttt{TestAnswer}: Câu trả lời trong bài kiểm tra
    \item \texttt{Class}: Lớp học
    \item \texttt{ClassMember}: Thành viên lớp học
    \item \texttt{ClassStudySet}: Liên kết studyset với class
    \item \texttt{ChatConversation}: Cuộc hội thoại với AI
    \item \texttt{ChatMessage}: Message trong conversation
    \item \texttt{Category}: Danh mục
    \item \texttt{AIGeneratedContents}: Nội dung được AI sinh ra
    \item \texttt{RefreshToken}: Refresh token cho authentication
    \item \texttt{ReattemptRequest}: Yêu cầu làm lại bài kiểm tra
\end{itemize}

Các quan hệ chính:
\begin{itemize}
    \item \texttt{User} kết hợp với \texttt{StudySet} (association, một-nhiều)
    \item \texttt{StudySet} hợp thành với \texttt{Term} (composition, một-nhiều)
    \item \texttt{User} kết hợp với \texttt{StudyActivity} (association, một-nhiều)
    \item \texttt{StudySet} kết hợp với \texttt{ProgressSummary} (association, một-nhiều)
    \item \texttt{StudySet} kết hợp với \texttt{Test} (association, một-nhiều)
    \item \texttt{User} kết hợp với \texttt{Class} thông qua \texttt{ClassMember} (association, nhiều-nhiều)
\end{itemize}

\section{Thiết kế chi tiết}
\label{section:4.2}

\subsection{Thiết kế giao diện}
\label{subsection:4.2.1}

Hệ thống Aikari được thiết kế để hỗ trợ các thiết bị và màn hình khác nhau, đảm bảo trải nghiệm người dùng tốt trên mọi nền tảng.

\subsubsection{Thông số kỹ thuật màn hình}

Ứng dụng được thiết kế để hỗ trợ:
\begin{itemize}
    \item \textbf{Độ phân giải}: Tối thiểu 1280x720 (HD), tối ưu cho 1920x1080 (Full HD) và cao hơn
    \item \textbf{Kích thước màn hình}: Hỗ trợ từ 13 inch (laptop) đến 27 inch (desktop)
    \item \textbf{Số lượng màu sắc}: Hỗ trợ đầy đủ 24-bit color (16.7 triệu màu)
    \item \textbf{Responsive design}: Tự động điều chỉnh layout cho mobile, tablet, và desktop
\end{itemize}

\subsubsection{Thống nhất thiết kế giao diện}

Hệ thống sử dụng Tailwind CSS v4 với OKLCH color space và tuân thủ các nguyên tắc thiết kế sau:

\textbf{Thiết kế nút và điều khiển}:
\begin{itemize}
    \item Nút chính (Primary): Sử dụng biến CSS \texttt{--color-primary}, bo góc \texttt{rounded-md} (6px), chiều cao mặc định \texttt{h-9} (36px), padding \texttt{px-4 py-2} (16px 8px)
    \item Nút phụ (Secondary): Sử dụng \texttt{--color-secondary}, cùng kích thước và bo góc như nút chính
    \item Nút nguy hiểm (Destructive): Sử dụng \texttt{--color-destructive}, cùng kích thước và bo góc
    \item Các kích thước nút: \texttt{sm} (h-8, px-3), \texttt{default} (h-9, px-4), \texttt{lg} (h-10, px-6)
    \item Input fields: Chiều cao \texttt{h-9}, bo góc \texttt{rounded-md} (6px), padding \texttt{px-3 py-1}, font size \texttt{text-base} (16px) trên mobile và \texttt{text-sm} (14px) trên desktop
\end{itemize}

\textbf{Vị trí hiển thị thông điệp phản hồi}:
\begin{itemize}
    \item Sử dụng thư viện Sonner với \texttt{position="top-right"}
    \item Thông báo thành công: Màu xanh lá (richColors)
    \item Thông báo lỗi: Màu đỏ (richColors)
    \item Thông báo cảnh báo: Màu vàng (richColors)
    \item Thông báo thông tin: Màu xanh dương (richColors)
\end{itemize}

\textbf{Phối màu (OKLCH color space)}:
\begin{itemize}
    \item Hệ thống sử dụng OKLCH color space để đảm bảo tính nhất quán và hỗ trợ dark mode
    \item Light mode: Primary (oklch(25.6\% 0.038 285.82)), Background (oklch(100\% 0 0)), Foreground (oklch(15.23\% 0.042 285.75))
    \item Dark mode: Primary (oklch(98.04\% 0.002 264.53)), Background (oklch(15.23\% 0.042 285.75)), Foreground (oklch(98.04\% 0.002 264.53))
    \item Destructive: Light mode (oklch(62.8\% 0.204 29.23)), Dark mode (oklch(42.09\% 0.153 29.7))
    \item Secondary: Light mode (oklch(96.84\% 0.004 264.53)), Dark mode (oklch(28.51\% 0.036 274.15))
    \item Border và Input: Sử dụng \texttt{--color-border} và \texttt{--color-input} tương ứng với theme
\end{itemize}

\textbf{Layout và spacing}:
\begin{itemize}
    \item Container: Sử dụng Tailwind \texttt{container mx-auto} (max-width responsive), các trang chính sử dụng \texttt{max-w-4xl} (896px) với padding \texttt{p-6} (24px)
    \item Grid system: Sử dụng CSS Grid và Flexbox thông qua Tailwind utilities
    \item Spacing: Sử dụng Tailwind spacing scale (4px base): 1 (4px), 2 (8px), 3 (12px), 4 (16px), 6 (24px), 8 (32px), 12 (48px), 16 (64px)
    \item Border radius: \texttt{rounded-md} (6px) cho các element chính, \texttt{--radius-lg} (8px) cho các element lớn, \texttt{--radius-sm} (4px) cho các element nhỏ
\end{itemize}

\textbf{Typography}:
\begin{itemize}
    \item Font chính: Inter (từ Google Fonts), weights từ 100 đến 900
    \item Font sizes: \texttt{text-xs} (12px), \texttt{text-sm} (14px), \texttt{text-base} (16px), \texttt{text-lg} (18px), \texttt{text-xl} (20px), \texttt{text-2xl} (24px), \texttt{text-4xl} (36px)
    \item Font weights: \texttt{font-normal} (400), \texttt{font-medium} (500), \texttt{font-bold} (700)
\end{itemize}

\subsubsection{Minh họa giao diện các chức năng chính}

\textbf{Trang đăng nhập}: Giao diện đăng nhập được thiết kế đơn giản, tập trung vào form đăng nhập với các trường username và password, nút đăng nhập, và liên kết quên mật khẩu. Layout căn giữa, sử dụng card với shadow nhẹ.

\textbf{Trang Dashboard}: Hiển thị tổng quan về tiến độ học tập, số lượng studysets, số lượng terms đã học, và thống kê học tập. Sử dụng các card để hiển thị thông tin, biểu đồ để minh họa tiến độ.

\textbf{Trang học tập (Study)}: Giao diện học flashcard với card hiển thị term, nút để xem đáp án, và các nút đánh giá (Again, Hard, Good, Easy) tương ứng với recall score. Giao diện tập trung, dễ sử dụng.

\textbf{Trang Chatbot}: Giao diện chat với AI, hiển thị lịch sử conversation, input box để nhập câu hỏi, và nút gửi. Messages được phân biệt giữa user và AI bằng màu sắc khác nhau.

\subsection{Thiết kế lớp}
\label{subsection:4.2.2}

Phần này trình bày thiết kế chi tiết các thuộc tính và phương thức cho các lớp quan trọng nhất của hệ thống.

\subsubsection{Lớp LearningService}

Lớp \texttt{LearningService} chịu trách nhiệm quản lý logic học tập và triển khai thuật toán SM-2.

\textbf{Thuộc tính chính}:
\begin{itemize}
    \item \texttt{db: Session}: Database session để truy xuất dữ liệu
\end{itemize}

\textbf{Phương thức chính}:
\begin{itemize}
    \item \texttt{record\_review(user\_id, term\_id, recall\_score)}: Ghi nhận một lần học, tính toán SM-2, cập nhật StudyActivity và ProgressSummary
    \item \texttt{get\_next\_term(user\_id, studyset\_id)}: Lấy term tiếp theo cần học (ưu tiên due terms)
    \item \texttt{calculate\_progress(user\_id, studyset\_id)}: Tính toán tiến độ học tập
    \item \texttt{get\_study\_statistics(user\_id, studyset\_id)}: Lấy thống kê học tập
\end{itemize}

\subsubsection{Lớp ChatbotService}

Lớp \texttt{ChatbotService} chịu trách nhiệm xử lý conversation với AI chatbot.

\textbf{Thuộc tính chính}:
\begin{itemize}
    \item \texttt{db: Session}: Database session
    \item \texttt{dify\_service: DifyService}: Service để gọi Dify API
\end{itemize}

\textbf{Phương thức chính}:
\begin{itemize}
    \item \texttt{process\_message(conversation\_id, message)}: Xử lý message từ user, xác định intent, và gọi Dify API
    \item \texttt{create\_conversation(user\_id, studyset\_id)}: Tạo conversation mới
    \item \texttt{get\_conversation\_history(conversation\_id)}: Lấy lịch sử conversation
    \item \texttt{update\_conversation\_state(conversation\_id, state, intent, params)}: Cập nhật state và intent của conversation
\end{itemize}

\subsubsection{Lớp SessionService}

Lớp \texttt{SessionService} chịu trách nhiệm quản lý learning sessions và session-based reviews.

\textbf{Thuộc tính chính}:
\begin{itemize}
    \item \texttt{session: Session}: Database session
\end{itemize}

\textbf{Phương thức chính}:
\begin{itemize}
    \item \texttt{start\_session(user\_id, studyset\_id, session\_type, total\_cards)}: Bắt đầu một learning session mới
    \item \texttt{record\_review(session\_id, term\_id, recall\_score)}: Ghi nhận review trong session
    \item \texttt{end\_session(session\_id)}: Kết thúc session và cập nhật tiến độ
    \item \texttt{get\_session\_statistics(session\_id)}: Lấy thống kê của session
\end{itemize}

\subsubsection{Biểu đồ trình tự cho use case "Học flashcard"}

Biểu đồ trình tự mô tả luồng xử lý khi người dùng học một flashcard được trình bày trong Hình \ref{fig:study-sequence}.

\begin{figure}[H]
    \centering
    \includegraphics[width=0.9\textwidth]{Hinhve/study-sequence.png}
    \caption{Biểu đồ trình tự cho use case "Học flashcard"}
    \label{fig:study-sequence}
\end{figure}

Luồng xử lý:
\begin{enumerate}
    \item User gửi request học flashcard từ frontend
    \item API Route nhận request, validate input
    \item Route gọi LearningService.get\_next\_term()
    \item LearningService truy vấn database để lấy term cần học
    \item Database trả về term
    \item LearningService trả về term cho Route
    \item Route trả về term cho frontend
    \item User đánh giá recall score
    \item Frontend gửi request với recall\_score
    \item Route gọi LearningService.record\_review()
    \item LearningService tính toán SM-2 (EF, interval, next\_review\_date)
    \item LearningService tạo StudyActivity mới
    \item LearningService cập nhật ProgressSummary
    \item LearningService trả về kết quả
    \item Route trả về response cho frontend
\end{enumerate}

\subsubsection{Biểu đồ trình tự cho use case "Chat với AI"}

Biểu đồ trình tự mô tả luồng xử lý khi người dùng chat với AI được trình bày trong Hình \ref{fig:chatbot-sequence}.

\begin{figure}[H]
    \centering
    \includegraphics[width=0.9\textwidth]{Hinhve/chatbot-sequence.png}
    \caption{Biểu đồ trình tự cho use case "Chat với AI"}
    \label{fig:chatbot-sequence}
\end{figure}

Luồng xử lý:
\begin{enumerate}
    \item User gửi message từ frontend
    \item API Route nhận request, validate input
    \item Route gọi ChatbotService.process\_message()
    \item ChatbotService lấy conversation từ database
    \item ChatbotService xác định intent và state
    \item ChatbotService gọi DifyService để xử lý message
    \item DifyService gửi request đến Dify API
    \item Dify API trả về response
    \item DifyService trả về response cho ChatbotService
    \item ChatbotService lưu message vào database
    \item ChatbotService cập nhật conversation state
    \item ChatbotService trả về response cho Route
    \item Route trả về response cho frontend
\end{enumerate}

\subsection{Thiết kế cơ sở dữ liệu}
\label{subsection:4.2.3}

\subsubsection{Biểu đồ thực thể liên kết (E-R Diagram)}

Biểu đồ thực thể liên kết mô tả các thực thể chính và quan hệ giữa chúng trong hệ thống được trình bày trong Hình \ref{fig:er-diagram}.

\begin{figure}[H]
    \centering
    \includegraphics[width=0.9\textwidth]{Hinhve/er-diagram.png}
    \caption{Biểu đồ thực thể liên kết của hệ thống Aikari}
    \label{fig:er-diagram}
\end{figure}

\textbf{Các thực thể chính}:

\begin{itemize}
    \item \textbf{User}: Người dùng hệ thống
    \begin{itemize}
        \item user\_id (PK), username, email, password, full\_name, role, phone\_numbers, address, city, country, preferences, created\_at, updated\_at
    \end{itemize}
    
    \item \textbf{StudySet}: Bộ flashcard
    \begin{itemize}
        \item studyset\_id (PK), title, description, owner\_id (FK), content\_type, category\_id (FK), created\_at, updated\_at
    \end{itemize}
    
    \item \textbf{Term}: Thẻ flashcard
    \begin{itemize}
        \item term\_id (PK), studyset\_id (FK), front, back, easiness\_factor, interval, repetitions, next\_review\_date, created\_at, updated\_at
    \end{itemize}
    
    \item \textbf{StudyActivity}: Ghi nhận mỗi lần học
    \begin{itemize}
        \item activity\_id (PK), user\_id (FK), term\_id (FK), studyset\_id (FK), recall\_score, easiness\_factor, interval, review\_date, created\_at
    \end{itemize}
    
    \item \textbf{ProgressSummary}: Tổng hợp tiến độ học tập
    \begin{itemize}
        \item progress\_id (PK), user\_id (FK), studyset\_id (FK), total\_terms, learned\_terms, mastered\_terms, last\_studied\_at, created\_at, updated\_at
    \end{itemize}
    
    \item \textbf{Test}: Bài kiểm tra
    \begin{itemize}
        \item test\_id (PK), studyset\_id (FK), title, description, test\_type, created\_by (FK), created\_at, updated\_at
    \end{itemize}
    
    \item \textbf{Class}: Lớp học
    \begin{itemize}
        \item class\_id (PK), class\_name, description, created\_by, owner\_user\_id (FK), is\_public, class\_code, created\_at, updated\_at
    \end{itemize}
    
    \item \textbf{ChatConversation}: Cuộc hội thoại với AI
    \begin{itemize}
        \item conversation\_id (PK), studyset\_id (FK), user\_id (FK), state, intent, collected\_params, created\_at, updated\_at
    \end{itemize}
    
    \item \textbf{LearningSession}: Learning session
    \begin{itemize}
        \item session\_id (PK), user\_id (FK), studyset\_id (FK), session\_type, status, total\_cards, reviewed\_cards, started\_at, ended\_at, created\_at, updated\_at
    \end{itemize}
    
    \item \textbf{SessionReview}: Review trong session
    \begin{itemize}
        \item review\_id (PK), session\_id (FK), term\_id (FK), recall\_score, review\_time, created\_at
    \end{itemize}
    
    \item \textbf{RefreshToken}: Refresh token
    \begin{itemize}
        \item token\_id (PK), user\_id (FK), token, expires\_at, created\_at
    \end{itemize}
\end{itemize}

\textbf{Các quan hệ chính}:

\begin{itemize}
    \item User --(1:N)--> StudySet: Một user có thể tạo nhiều studysets
    \item StudySet --(1:N)--> Term: Một studyset có nhiều terms
    \item User --(1:N)--> StudyActivity: Một user có nhiều study activities
    \item Term --(1:N)--> StudyActivity: Một term có nhiều study activities
    \item User --(1:N)--> ProgressSummary: Một user có nhiều progress summaries
    \item StudySet --(1:N)--> ProgressSummary: Một studyset có nhiều progress summaries
    \item StudySet --(1:N)--> Test: Một studyset có nhiều tests
    \item User --(N:M)--> Class: Nhiều users tham gia nhiều classes (thông qua ClassMember)
    \item StudySet --(N:M)--> Class: Nhiều studysets được gán cho nhiều classes (thông qua ClassStudySet)
    \item User --(1:N)--> ChatConversation: Một user có nhiều conversations
    \item StudySet --(1:N)--> ChatConversation: Một studyset có nhiều conversations
    \item User --(1:N)--> LearningSession: Một user có nhiều learning sessions
    \item StudySet --(1:N)--> LearningSession: Một studyset có nhiều learning sessions
    \item LearningSession --(1:N)--> SessionReview: Một session có nhiều reviews
    \item Term --(1:N)--> SessionReview: Một term có nhiều session reviews
    \item User --(1:N)--> RefreshToken: Một user có nhiều refresh tokens
    \item Test --(1:N)--> TestQuestion: Một test có nhiều questions
    \item Test --(1:N)--> TestAttempt: Một test có nhiều attempts
    \item TestAttempt --(1:N)--> TestAnswer: Một attempt có nhiều answers
    \item TestQuestion --(1:N)--> TestAnswer: Một question có nhiều answers
    \item User --(1:N)--> ReattemptRequest: Một user có nhiều reattempt requests
    \item Test --(1:N)--> ReattemptRequest: Một test có nhiều reattempt requests
\end{itemize}

\subsubsection{Thiết kế cơ sở dữ liệu PostgreSQL}

Hệ thống sử dụng PostgreSQL làm hệ quản trị cơ sở dữ liệu. Các bảng được thiết kế dựa trên biểu đồ E-R và được quản lý thông qua SQLModel ORM và Alembic migrations.

\textbf{Các bảng chính}:

\begin{itemize}
    \item \texttt{User}: Lưu trữ thông tin người dùng
    \item \texttt{StudySet}: Lưu trữ thông tin bộ flashcard
    \item \texttt{Term}: Lưu trữ thông tin thẻ flashcard
    \item \texttt{StudyActivity}: Lưu trữ lịch sử học tập
    \item \texttt{ProgressSummary}: Lưu trữ tổng hợp tiến độ
    \item \texttt{LearningSession}: Lưu trữ thông tin learning sessions
    \item \texttt{SessionReview}: Lưu trữ reviews trong session
    \item \texttt{Test}: Lưu trữ thông tin bài kiểm tra
    \item \texttt{TestQuestion}: Lưu trữ câu hỏi trong bài kiểm tra
    \item \texttt{TestAttempt}: Lưu trữ lần làm bài kiểm tra
    \item \texttt{TestAnswer}: Lưu trữ câu trả lời trong bài kiểm tra
    \item \texttt{ReattemptRequest}: Lưu trữ yêu cầu làm lại bài kiểm tra
    \item \texttt{Class}: Lưu trữ thông tin lớp học
    \item \texttt{ClassMember}: Lưu trữ thành viên lớp học
    \item \texttt{ClassStudySet}: Liên kết studyset với class
    \item \texttt{ChatConversation}: Lưu trữ conversation với AI
    \item \texttt{ChatMessage}: Lưu trữ messages trong conversation
    \item \texttt{Category}: Lưu trữ danh mục
    \item \texttt{AIGeneratedContents}: Lưu trữ nội dung được AI sinh ra
    \item \texttt{RefreshToken}: Lưu trữ refresh tokens cho authentication
\end{itemize}

\textbf{Indexes và constraints}:

\begin{itemize}
    \item Primary keys: Tất cả các bảng sử dụng UUID làm primary key
    \item Foreign keys: Được định nghĩa rõ ràng để đảm bảo tính toàn vẹn dữ liệu
    \item Indexes: Được tạo trên các cột thường xuyên truy vấn (user\_id, studyset\_id, email, username, next\_review\_date)
    \item Unique constraints: email, username, class\_code
    \item Check constraints: Đảm bảo giá trị hợp lệ cho các enum fields
\end{itemize}

\section{Xây dựng ứng dụng}
\label{section:4.3}

\subsection{Thư viện và công cụ sử dụng}
\label{subsection:4.3.1}

Hệ thống Aikari sử dụng các công cụ, ngôn ngữ lập trình, API, thư viện, IDE, và công cụ kiểm thử như được trình bày trong Bảng \ref{table:tools}.

\begin{table}[H]
\centering
\begin{tabular}{lll}
\hline
\textbf{Mục đích} & \textbf{Công cụ} & \textbf{Phiên bản} \\ \hline
Ngôn ngữ lập trình (Backend) & Python & 3.10+ \\
Ngôn ngữ lập trình (Frontend) & TypeScript & 5.x \\
Web Framework (Backend) & FastAPI & 0.114.2+ \\
UI Framework (Frontend) & React Router & v7 \\
UI Library & React & 19.x \\
ORM & SQLModel & Latest \\
Database & PostgreSQL & 12+ \\
Database Migrations & Alembic & Latest \\
State Management & Redux Toolkit & Latest \\
Styling & Tailwind CSS & v4 \\
UI Components & Radix UI & Latest \\
Build Tool (Frontend) & Vite & Latest \\
Package Manager (Python) & uv & Latest \\
Package Manager (Node) & npm & Latest \\
Testing Framework (Backend) & pytest & Latest \\
Testing Framework (Frontend) & Vitest & Latest \\
Code Formatter (Python) & ruff & Latest \\
Linter (Python) & ruff & Latest \\
Type Checker (Python) & mypy & Latest \\
Type Checker (TypeScript) & TypeScript & 5.x \\
AI Platform & Dify & Latest \\
Authentication & JWT & Latest \\
Password Hashing & passlib & Latest \\
Email Templates & MJML & Latest \\ \hline
\end{tabular}
\caption{Danh sách thư viện và công cụ sử dụng}
\label{table:tools}
\end{table}

\subsection{Kết quả đạt được}
\label{subsection:4.3.2}

Sau khi hoàn thiện hệ thống, người dùng có thể thực hiện các chức năng sau:

Ứng dụng Aikari là một \textbf{Hệ thống Học tập Flashcard tích hợp Trợ lý ảo AI}, nơi người dùng có thể tạo, quản lý và học tập hiệu quả thông qua flashcard với thuật toán Spaced Repetition (SM-2). Hệ thống hỗ trợ tìm kiếm và duyệt các bộ flashcard theo nhiều tiêu chí khác nhau như danh mục (category), chủ đề, và người tạo. Hệ thống sẽ hiển thị danh sách các bộ flashcard phù hợp và cung cấp đầy đủ thông tin chi tiết về từng bộ flashcard, bao gồm số lượng terms, tiến độ học tập, thống kê nhớ/quên, và các thông tin liên quan khác.

Người học (Student) có thể tạo và quản lý các bộ flashcard cá nhân, tổ chức chúng theo danh mục, và học tập theo nhiều chế độ khác nhau (học trực tiếp, quick review, study mode). Khi học flashcard, người dùng có thể đánh giá mức độ nhớ của mình bằng các mức: Again (1), Hard (2), Good (3), Easy (4). Hệ thống tự động tính toán và cập nhật next\_review\_date dựa trên thuật toán SM-2, đảm bảo việc ôn tập được tối ưu hóa theo khả năng ghi nhớ của từng người dùng. Người học có thể xem thống kê tiến độ học tập cá nhân, bao gồm số flashcard đã học, tỉ lệ nhớ/quên, và thời gian học.

Trong quá trình học, người dùng có thể tương tác với \textbf{trợ lý ảo AI} được tích hợp trực tiếp trong màn hình học flashcard. Chatbot AI có thể giúp giải thích thuật ngữ, trả lời câu hỏi về nội dung đang học, và sinh nội dung học tập như bài test, quiz, hoặc đoạn văn (paragraph) dựa trên flashcard đang học. Người dùng có thể gửi các yêu cầu đặc biệt cho AI, chẳng hạn như "Tạo test cho studyset này" hoặc "Giải thích thuật ngữ X"; những yêu cầu này sẽ được xử lý thông qua Dify workflows và trả về kết quả phù hợp.

Giáo viên (Teacher) có toàn bộ chức năng của Student, đồng thời có thêm khả năng tạo và quản lý lớp học. Giáo viên có thể tạo lớp học, thêm studysets vào lớp, quản lý danh sách học viên, và xem thống kê tiến độ học tập của toàn bộ lớp. Hệ thống cung cấp các thống kê chi tiết như tỉ lệ nhớ/quên của từng học viên, số flashcard đã hoàn thành, và mức độ tiến bộ theo thời gian. Học sinh có thể tìm kiếm và tham gia lớp học bằng class code, sau đó xem và học các studysets được giáo viên gán cho lớp.

Hệ thống cũng hỗ trợ tạo và làm bài kiểm tra (Test), cho phép giáo viên tạo các bài test với nhiều loại câu hỏi khác nhau (multiple choice, true/false, v.v.) dựa trên các flashcard trong studyset. Học sinh có thể làm bài kiểm tra, xem kết quả, và yêu cầu làm lại bài nếu cần. Các kết quả và lịch sử làm bài được lưu trữ và hiển thị, giúp người dùng theo dõi tiến độ học tập của mình.

Vai trò quản trị cao nhất thuộc về Admin, người có quyền kiểm soát toàn bộ hệ thống, bao gồm giám sát hoạt động, quản lý người dùng, phân quyền truy cập, và xem thống kê tổng quan về hệ thống như số lượng người dùng, hoạt động học tập, và hiệu suất sử dụng hệ thống.

Việc phân tách rõ ràng các chức năng theo từng vai trò (Student, Teacher, Admin) giúp hệ thống vận hành hiệu quả, đồng thời đảm bảo tính bảo mật và trải nghiệm người dùng tối ưu. Hệ thống được xây dựng với kiến trúc module hóa, dễ mở rộng và bảo trì, sẵn sàng cho triển khai thực tế.

Thống kê về ứng dụng được trình bày trong Bảng \ref{table:statistics}.

\begin{table}[H]
\centering
\begin{tabular}{ll}
\hline
\textbf{Thông tin} & \textbf{Giá trị} \\ \hline
Số dòng code (Backend) & ~15,000+ dòng \\
Số dòng code (Frontend) & ~20,000+ dòng \\
Số lớp (Backend) & 30+ lớp \\
Số component (Frontend) & 50+ components \\
Số gói (Backend) & 10+ gói \\
Số gói (Frontend) & 8+ gói \\
Số lượng API endpoints & 40+ endpoints \\ \hline
\end{tabular}
\caption{Thống kê về ứng dụng}
\label{table:statistics}
\end{table}

\subsection{Minh họa các chức năng chính}
\label{subsection:4.3.3}

\subsubsection{Chức năng đăng nhập và xác thực}

Giao diện đăng nhập cho phép người dùng đăng nhập vào hệ thống bằng username và password. Sau khi đăng nhập thành công, hệ thống cấp JWT token để xác thực các request tiếp theo. Giao diện được thiết kế đơn giản, dễ sử dụng với validation rõ ràng cho các trường input.

\subsubsection{Chức năng quản lý StudySet}

Giao diện quản lý StudySet cho phép người dùng tạo, chỉnh sửa, xóa, và xem danh sách các bộ flashcard. Người dùng có thể thêm terms vào studyset, chỉnh sửa nội dung, và xóa các terms không cần thiết. Giao diện sử dụng table để hiển thị danh sách, form để tạo/chỉnh sửa, và modal để xác nhận xóa.

\subsubsection{Chức năng học tập với SM-2}

Giao diện học tập hiển thị flashcard với mặt trước (front) và mặt sau (back). Người dùng có thể xem đáp án và đánh giá mức độ nhớ của mình bằng các nút: Again (1), Hard (2), Good (3), Easy (4). Hệ thống tự động tính toán và cập nhật next\_review\_date dựa trên thuật toán SM-2. Giao diện hiển thị tiến độ học tập và số lượng terms còn lại.

\subsubsection{Chức năng Chatbot AI}

Giao diện chatbot cho phép người dùng tương tác với AI để được hỗ trợ học tập. Chatbot có thể giúp giải thích thuật ngữ, tạo câu hỏi, và cung cấp thông tin về studyset. Giao diện hiển thị lịch sử conversation, input box để nhập câu hỏi, và nút gửi. Messages được phân biệt giữa user và AI bằng màu sắc và vị trí.

\subsubsection{Chức năng quản lý lớp học}

Giao diện quản lý lớp học cho phép giáo viên tạo lớp, thêm studysets vào lớp, và quản lý thành viên. Học sinh có thể tham gia lớp bằng class code và xem các studysets được gán cho lớp. Giao diện hiển thị danh sách lớp, thông tin chi tiết lớp, và danh sách thành viên.

\section{Kiểm thử}
\label{section:4.4}

Hệ thống Aikari sử dụng các kỹ thuật kiểm thử khác nhau để đảm bảo chất lượng và độ tin cậy của code. Phần này trình bày về các kỹ thuật kiểm thử đã được áp dụng và kết quả kiểm thử cho các chức năng quan trọng nhất.

\subsection{Kỹ thuật kiểm thử sử dụng}

Hệ thống sử dụng các kỹ thuật kiểm thử sau:

\begin{enumerate}
    \item \textbf{Unit Testing}: Kiểm thử từng đơn vị code (function, method) độc lập, sử dụng \textbf{pytest} framework.
    
    \item \textbf{Integration Testing}: Kiểm thử tích hợp giữa các components, đặc biệt là API endpoints với database và services.
    
    \item \textbf{API Testing}: Kiểm thử các REST API endpoints sử dụng \textbf{TestClient} của FastAPI.
    
    \item \textbf{Test Fixtures}: Sử dụng pytest fixtures để setup test data, database sessions, và authentication tokens.
\end{enumerate}

\subsection{Thiết kế test cases cho các chức năng quan trọng}

Hệ thống đã được phát triển xong các chức năng cơ bản như Authentication, User Management, SM-2 algorithm, learning sessions, và chatbot. Phần này trình bày về các test cases đã được thiết kế và thực hiện cho các chức năng quan trọng nhất.

\subsubsection{Chức năng 1: Authentication và Authorization}

Authentication là chức năng cơ bản và quan trọng nhất của hệ thống, đảm bảo chỉ người dùng hợp lệ mới có thể truy cập hệ thống.

\textbf{Test Case 1.1: Đăng nhập thành công}
\begin{itemize}
    \item \textbf{Mục đích}: Kiểm tra người dùng có thể đăng nhập với credentials hợp lệ
    \item \textbf{Input}: Username và password hợp lệ
    \item \textbf{Expected Output}: HTTP 200, trả về access\_token
    \item \textbf{Kỹ thuật}: API Testing với TestClient
\end{itemize}

\textbf{Test Case 1.2: Đăng nhập với mật khẩu sai}
\begin{itemize}
    \item \textbf{Mục đích}: Kiểm tra hệ thống từ chối đăng nhập với mật khẩu sai
    \item \textbf{Input}: Username hợp lệ, password sai
    \item \textbf{Expected Output}: HTTP 400, không trả về token
    \item \textbf{Kỹ thuật}: API Testing, Negative Testing
\end{itemize}

\textbf{Test Case 1.3: Sử dụng access token}
\begin{itemize}
    \item \textbf{Mục đích}: Kiểm tra token có thể được sử dụng để xác thực
    \item \textbf{Input}: Access token hợp lệ trong Authorization header
    \item \textbf{Expected Output}: HTTP 200, trả về thông tin user
    \item \textbf{Kỹ thuật}: API Testing với authentication headers
\end{itemize}

\textbf{Test Case 1.4: Khôi phục mật khẩu}
\begin{itemize}
    \item \textbf{Mục đích}: Kiểm tra chức năng khôi phục mật khẩu
    \item \textbf{Input}: Email hợp lệ
    \item \textbf{Expected Output}: HTTP 200, email được gửi (mocked)
    \item \textbf{Kỹ thuật}: API Testing với mocking email service
\end{itemize}

\textbf{Test Case 1.5: Đặt lại mật khẩu}
\begin{itemize}
    \item \textbf{Mục đích}: Kiểm tra người dùng có thể đặt lại mật khẩu với token hợp lệ
    \item \textbf{Input}: Reset token hợp lệ và mật khẩu mới
    \item \textbf{Expected Output}: HTTP 200, mật khẩu được cập nhật trong database
    \item \textbf{Kỹ thuật}: Integration Testing với database verification
\end{itemize}

\textbf{Test Case 1.6: Đặt lại mật khẩu với token không hợp lệ}
\begin{itemize}
    \item \textbf{Mục đích}: Kiểm tra hệ thống từ chối đặt lại mật khẩu với token không hợp lệ
    \item \textbf{Input}: Token không hợp lệ
    \item \textbf{Expected Output}: HTTP 400, "Invalid token"
    \item \textbf{Kỹ thuật}: Negative Testing
\end{itemize}

\subsubsection{Chức năng 2: Quản lý người dùng}

Quản lý người dùng bao gồm CRUD operations và phân quyền, đảm bảo chỉ admin mới có thể quản lý users.

\textbf{Test Case 2.1: Tạo user mới (Admin)}
\begin{itemize}
    \item \textbf{Mục đích}: Kiểm tra admin có thể tạo user mới
    \item \textbf{Input}: User data hợp lệ, admin token
    \item \textbf{Expected Output}: HTTP 200, user được tạo trong database
    \item \textbf{Kỹ thuật}: Integration Testing với database verification
\end{itemize}

\textbf{Test Case 2.2: Tạo user với email đã tồn tại}
\begin{itemize}
    \item \textbf{Mục đích}: Kiểm tra hệ thống từ chối tạo user trùng email
    \item \textbf{Input}: Email đã tồn tại
    \item \textbf{Expected Output}: HTTP 400, error message
    \item \textbf{Kỹ thuật}: Negative Testing
\end{itemize}

\textbf{Test Case 2.3: User thường không thể tạo user}
\begin{itemize}
    \item \textbf{Mục đích}: Kiểm tra phân quyền, chỉ admin mới tạo được user
    \item \textbf{Input}: User data, normal user token
    \item \textbf{Expected Output}: HTTP 403, "The user doesn't have enough privileges"
    \item \textbf{Kỹ thuật}: Authorization Testing
\end{itemize}

\textbf{Test Case 2.4: Cập nhật thông tin user}
\begin{itemize}
    \item \textbf{Mục đích}: Kiểm tra user có thể cập nhật thông tin của mình
    \item \textbf{Input}: Updated user data, user token
    \item \textbf{Expected Output}: HTTP 200, thông tin được cập nhật trong database
    \item \textbf{Kỹ thuật}: Integration Testing
\end{itemize}

\textbf{Test Case 2.5: Đổi mật khẩu}
\begin{itemize}
    \item \textbf{Mục đích}: Kiểm tra user có thể đổi mật khẩu
    \item \textbf{Input}: Current password và new password
    \item \textbf{Expected Output}: HTTP 200, mật khẩu được hash và lưu vào database
    \item \textbf{Kỹ thuật}: Integration Testing với password verification
\end{itemize}

\textbf{Test Case 2.6: Xóa user}
\begin{itemize}
    \item \textbf{Mục đích}: Kiểm tra user có thể xóa tài khoản của mình
    \item \textbf{Input}: User token
    \item \textbf{Expected Output}: HTTP 200, user bị xóa khỏi database
    \item \textbf{Kỹ thuật}: Integration Testing với database verification
\end{itemize}

\textbf{Test Case 2.7: Đăng ký user mới}
\begin{itemize}
    \item \textbf{Mục đích}: Kiểm tra người dùng có thể đăng ký tài khoản mới
    \item \textbf{Input}: Email, password, full\_name hợp lệ
    \item \textbf{Expected Output}: HTTP 200, user được tạo trong database
    \item \textbf{Kỹ thuật}: Integration Testing
\end{itemize}


\subsubsection{Chức năng 3: Thuật toán SM-2 - Test Cases chi tiết}

Các test cases cho thuật toán SM-2 được thiết kế dựa trên logic thực tế của \texttt{SpacedRepetitionSM2.calculate\_next\_review()} và \texttt{LearningService.record\_review()}. Các test cases này đã được implement và kiểm thử kỹ lưỡng để đảm bảo tính chính xác của thuật toán.

\textbf{Test Case 3.1: Tính toán Easiness Factor và Interval với các recall score khác nhau}
\begin{itemize}
    \item \textbf{Mục đích}: Kiểm tra công thức tính EF và interval đúng với các recall score khác nhau, bao gồm cả trường hợp reset khi quên
    \item \textbf{Input}: 
    \begin{itemize}
        \item Trường hợp 1: recall\_score = 5, current\_ef = 2.5, current\_interval = 1, repetitions = 0
        \item Trường hợp 2: recall\_score = 4, current\_ef = 2.5, current\_interval = 6, repetitions = 2
        \item Trường hợp 3: recall\_score = 2, current\_ef = 2.5, current\_interval = 10, repetitions = 5
    \end{itemize}
    \item \textbf{Logic kiểm thử}: 
    \begin{itemize}
        \item Trường hợp 1: new\_ef = 2.6 (tăng 0.1), new\_repetitions = 1, new\_interval = 1 ngày
        \item Trường hợp 2: new\_ef = 2.5, new\_repetitions = 3, new\_interval = int(6 * 2.5) = 15 ngày
        \item Trường hợp 3: new\_ef = 2.18 (giảm), new\_interval = 1, new\_repetitions = 0 (reset vì recall\_score < 3)
    \end{itemize}
    \item \textbf{Expected Output}: Các giá trị EF, interval, repetitions được tính đúng theo công thức SM-2, next\_review\_date được tính dựa trên UTC+7 local timezone
    \item \textbf{Kỹ thuật}: Unit Testing cho \texttt{SpacedRepetitionSM2.calculate\_next\_review()}
\end{itemize}

\textbf{Test Case 3.2: Ghi nhận review và cập nhật progress}
\begin{itemize}
    \item \textbf{Mục đích}: Kiểm tra \texttt{LearningService.record\_review()} tạo StudyActivity đúng và cập nhật ProgressSummary với thống kê chính xác
    \item \textbf{Input}: user\_id, studyset\_id, term\_id, recall\_score = 4, hint\_used = False, response\_time = 2.5
    \item \textbf{Logic kiểm thử}: 
    \begin{itemize}
        \item Lấy previous activity (nếu có) để lấy current\_ef, current\_interval, repetitions
        \item Gọi \texttt{SpacedRepetitionSM2.calculate\_next\_review()} để tính toán
        \item Tạo StudyActivity mới với các giá trị: ef, interval, repetitions, next\_review\_date, recall\_score, response\_time
        \item Lưu vào database và commit
        \item Gọi \texttt{update\_progress\_summary()} để cập nhật: mastered\_terms, reviewing\_terms, forgotten\_terms, average\_recall\_score, completion\_rate
    \end{itemize}
    \item \textbf{Expected Output}: StudyActivity được tạo với tất cả fields đúng, ProgressSummary được cập nhật với các giá trị thống kê chính xác
    \item \textbf{Kỹ thuật}: Integration Testing với database verification
\end{itemize}

\textbf{Test Case 3.3: Lấy term tiếp theo cần học - ưu tiên due terms}
\begin{itemize}
    \item \textbf{Mục đích}: Kiểm tra \texttt{LearningService.get\_next\_term\_to\_review()} chọn đúng term theo thứ tự ưu tiên
    \item \textbf{Input}: User có nhiều terms với next\_review\_date khác nhau, một số chưa học bao giờ
    \item \textbf{Logic kiểm thử}: 
    \begin{itemize}
        \item Priority 1: Terms có next\_review\_date <= today (UTC+7) - chọn term có next\_review\_date sớm nhất
        \item Priority 2: Terms chưa học bao giờ (không có StudyActivity)
        \item Priority 3: Terms đã học nhưng chưa đến hạn - chọn term có EF thấp nhất (hardest)
    \end{itemize}
    \item \textbf{Expected Output}: Term được chọn đúng theo thứ tự ưu tiên, xử lý timezone UTC+7 đúng
    \item \textbf{Kỹ thuật}: Integration Testing với database
\end{itemize}

\subsubsection{Chức năng 4: Dify AI Integration và Chatbot Flow - Test Cases chi tiết}

Các test cases cho Dify AI Integration và Chatbot Flow được thiết kế dựa trên logic thực tế của \texttt{DifyService}, \texttt{ChatbotService}, và \texttt{GenerationService}. Các test cases này đã được implement để đảm bảo tích hợp Dify API hoạt động đúng và conversation flow được xử lý chính xác.

\textbf{Test Case 4.1: Chat completion với Dify API thành công}
\begin{itemize}
    \item \textbf{Mục đích}: Kiểm tra \texttt{DifyService.chat\_completion()} gửi request đúng format và nhận response từ Dify API, bao gồm cả error handling
    \item \textbf{Input}: 
    \begin{itemize}
        \item Trường hợp thành công: query = "Giải thích thuật ngữ photosynthesis", user\_id = "user123", conversation\_id = None
        \item Trường hợp lỗi: Request với invalid API key hoặc network error
    \end{itemize}
    \item \textbf{Logic kiểm thử}: 
    \begin{itemize}
        \item Tạo request body với format: \{"inputs": \{\}, "query": query, "response\_mode": "blocking", "user": user\_id\}
        \item Gửi POST request đến endpoint "chat-messages" với Authorization header (Bearer token)
        \item Parse response JSON từ Dify API
        \item Xử lý lỗi: catch \texttt{httpx.HTTPStatusError} và \texttt{httpx.RequestError}, log error và raise exception
    \end{itemize}
    \item \textbf{Expected Output}: 
    \begin{itemize}
        \item Trường hợp thành công: Response chứa "answer", "conversation\_id", "id" (message\_id), "metadata"
        \item Trường hợp lỗi: Exception được raise với thông báo lỗi rõ ràng, không crash hệ thống
    \end{itemize}
    \item \textbf{Kỹ thuật}: Integration Testing với mocked Dify API hoặc test Dify API thực tế, mocked HTTP errors
\end{itemize}

\textbf{Test Case 4.2: Workflow run để sinh nội dung (test và paragraph)}
\begin{itemize}
    \item \textbf{Mục đích}: Kiểm tra \texttt{DifyService.run\_workflow()} chạy workflow để sinh test và paragraph từ studyset
    \item \textbf{Input}: 
    \begin{itemize}
        \item Sinh test: inputs = \{"studyset\_id": "ss123", "total\_questions": 10, "question\_types": ["MULTIPLE\_CHOICE"], "time\_limit": 20\}
        \item Sinh paragraph: inputs = \{"studyset\_id": "ss123", "term\_id": "term456"\}
    \end{itemize}
    \item \textbf{Logic kiểm thử}: 
    \begin{itemize}
        \item Tạo request body: \{"inputs": inputs, "response\_mode": "blocking", "user": user\_id\}
        \item Gửi POST request đến endpoint "workflows/run"
        \item Parse response chứa generated content (test với questions/answers hoặc paragraph text)
    \end{itemize}
    \item \textbf{Expected Output}: 
    \begin{itemize}
        \item Sinh test: Response chứa test content với format JSON, bao gồm questions, answers, metadata
        \item Sinh paragraph: Response chứa paragraph text dài, có context từ studyset và term
    \end{itemize}
    \item \textbf{Kỹ thuật}: Integration Testing với mocked Dify workflow
\end{itemize}

\textbf{Test Case 4.3: Tạo và quản lý conversation, lưu messages vào database}
\begin{itemize}
    \item \textbf{Mục đích}: Kiểm tra toàn bộ flow tạo conversation, xử lý messages, cập nhật state, và lưu vào database
    \item \textbf{Input}: studyset\_id, user\_id, message = "Tạo test", conversation\_id = None
    \item \textbf{Logic kiểm thử}: 
    \begin{itemize}
        \item Tạo conversation mới: \texttt{ChatbotService.get\_or\_create\_conversation()} tạo ChatConversation với state = INITIAL
        \item Detect intent: \texttt{detect\_user\_intent()} từ message keywords
        \item Lưu user message: Tạo ChatMessage với role = "user", content = message
        \item Xử lý message: Update conversation state theo flow (INITIAL → WAITING\_INTENT → COLLECTING\_TEST\_PARAMS → GENERATING → COMPLETED)
        \item Gọi Dify API: Chat completion hoặc workflow run tùy intent
        \item Lưu AI response: Tạo ChatMessage với role = "assistant", content = answer từ Dify
        \item Commit tất cả vào database
    \end{itemize}
    \item \textbf{Expected Output}: 
    \begin{itemize}
        \item ChatConversation được tạo/lấy đúng với state và intent được cập nhật
        \item Cả user message và AI response được lưu vào ChatMessage table
        \item Conversation state chuyển đổi đúng theo flow
    \end{itemize}
    \item \textbf{Kỹ thuật}: Integration Testing với database verification, mocked Dify API
\end{itemize}

Các test cases này đã được implement và kiểm thử kỹ lưỡng, đảm bảo tính chính xác của thuật toán SM-2 và tích hợp Dify AI. Tất cả test cases đều PASS, chứng minh hệ thống hoạt động đúng theo thiết kế.

\end{document}
